\chapter{What can be done with one qubit}\label{ch:applications1}

\section{Photonic qubits: polarization}\label{sec:photonic}
Light is quantized; its particles are called photons. A photon in a given spatial mode has two polarization states (horizontal and vertical
linear polarizations, as in 3D glasses; the circular basis can be used as well), so each photon is a two-level system, a qubit. Quantum gates
acting on photons are implemented by waveplates (quarter- and half-waveplates), made of birefringent materials, which have different
refractive indices for the two polarizations and therefore accumulate different phases. Projective measurements are implemented by
polarizers and polarizing beam splitters.

\subsection{Jones vectors and the Poincaré sphere}
\begin{origin}[: from Maxwell's equations to a qubit]
A monochromatic plane wave propagating along $z$ has the electric field
$\vec E(z,t)=\operatorname{Re}\left[(E_H\vec e_x+E_V\vec e_y)\e^{\ii(kz-\omega t)}\right]$. The complex two-component \emph{Jones vector}
$(E_H,E_V)$, normalized to unit intensity, describes the polarization. Its global phase is irrelevant to the polarization (it is just a
shift of the origin of time). The polarization of a classical wave is therefore described by exactly the same mathematics as a qubit:
a normalized vector of $\C^2$ modulo a global phase. For a single photon, the Jones vector becomes the quantum state
$\ket\psi=\alpha\ket H+\beta\ket V$ and $|\alpha|^2$ becomes a probability.
\end{origin}
We identify $\ket H=\ket0$, $\ket V=\ket1$. Then
\begin{center}
\begin{tabular}{@{}lll@{}}\toprule
Polarization & State & Point on the Poincaré sphere\\ \midrule
horizontal, vertical & $\ket H$, $\ket V$ & $\pm\vec e_z$\\
diagonal, antidiagonal ($\pm45^\circ$) & $\ket{D},\ket{A}=(\ket H\pm\ket V)/\sqrt2$ & $\pm\vec e_x$\\
right, left circular & $\ket{R},\ket L=(\ket H\pm\ii\ket V)/\sqrt2$ & $\pm\vec e_y$\\ \bottomrule
\end{tabular}
\end{center}
(The handedness convention for circular polarization varies between textbooks.) The components of the Bloch vector are the normalized
\emph{Stokes parameters}, measured classically with polarizers: $S_1=I_H-I_V$, $S_2=I_D-I_A$, $S_3=I_R-I_L$, divided by the total intensity.
A linear polarization at angle $\eta$ from the horizontal, $\cos\eta\ket H+\sin\eta\ket V$, corresponds by \cref{eq:blochstate} to
$\theta=2\eta$, $\varphi=0$: a point on the equator-crossing great circle in the $xz$ plane at polar angle $2\eta$.

The pseudospin is thus rotated by a \emph{double angle} with respect to the polarization: orthogonal polarizations ($H$ and $V$) differ by
$90^\circ$ in real space, while orthogonal qubit states are $180^\circ$ apart on the sphere. One needs to get used to it when working with photons.

\subsection{Waveplates as quantum gates}
A plate of thickness $L$ of a birefringent crystal with indices $n_f$ (fast axis) and $n_s$ (slow axis) delays the slow component by the phase
\begin{equation}
  \delta=\frac{2\pi(n_s-n_f)L}{\lambda}.
\end{equation}
With the fast axis horizontal, the plate acts as $J_0(\delta)=\diag(1,\e^{\ii\delta})=\e^{\ii\delta/2}R_z(\delta)$: a rotation of the Poincaré
sphere about the $z$ axis by $\delta$. A half-wave plate (HWP) has $\delta=\pi$ ($J_0=Z$), a quarter-wave plate (QWP) has $\delta=\pi/2$ ($J_0=S$).

A plate rotated by an angle $\alpha$ in the laboratory is $J(\alpha)=\mathcal R(\alpha)J_0\mathcal R(-\alpha)$ with the real rotation matrix
$\mathcal R(\alpha)=\begin{pmatrix}\cos\alpha&-\sin\alpha\\ \sin\alpha&\cos\alpha\end{pmatrix}=R_y(2\alpha)$: a physical rotation by $\alpha$ is a rotation of the
sphere by $2\alpha$ about $y$. For a half-wave plate,
\begin{equation}
  J_{\rm HWP}(\alpha)=\begin{pmatrix}\cos2\alpha&\sin2\alpha\\ \sin2\alpha&-\cos2\alpha\end{pmatrix}.
\end{equation}
It maps $\ket H$ to the linear polarization at $2\alpha$: a HWP at angle $\alpha$ \emph{mirrors} the polarization about its axis. Special cases:
at $\alpha=0$ it is the $Z$ gate; at $\alpha=45^\circ$ it is the $X$ gate ($H\leftrightarrow V$); at $\alpha=22.5^\circ$ it is exactly the
\emph{Hadamard gate}. A QWP at $45^\circ$ transforms linear into circular polarization. With a QWP--HWP--QWP sequence, any $SU(2)$ rotation
can be realised (this is the optical version of the Euler decomposition, \cref{thm:euler}).

\subsection{Measurement: Malus's law}
A polarizer along the direction $\eta_0$ realizes the projective measurement onto $\ket{\eta_0}=\cos\eta_0\ket H+\sin\eta_0\ket V$. The
equivalent of \cref{eq:projector} is \emph{Malus's law}: the probability of transmission is
\begin{equation}
  p=|\braket{\eta_0}\psi|^2=\cos^2\eta=\cos^2\frac\Theta2,
\end{equation}
where $\eta$ is the angle between the polarization of light and the transmission axis of the polarizer, and $\Theta=2\eta$ is the angle between
the pseudospin of the light and the pseudospin eigenvector of the polarizer. For a classical beam, Malus's law (1809) gives the transmitted
intensity; for a single photon it gives a probability.

\section{Random number generation}
\subsection{Why true randomness matters}
Random numbers are an essential resource in many areas: cryptographic keys and nonces, quantum key distribution, Monte Carlo simulations,
randomized algorithms, statistical sampling and lotteries. Most random numbers used in practice come from \emph{pseudo-random number generators}.
These are deterministic algorithms that expand a short seed into a long sequence that looks random. Anyone who knows the algorithm and the
seed can reproduce the whole sequence, which is a serious weakness in cryptography, and the sequence is eventually periodic.

Classical physical generators, based for instance on thermal noise or chaotic circuits, avoid the explicit seed. Their unpredictability,
however, comes only from our ignorance of the microscopic details: the underlying dynamics is deterministic in principle, and it is hard to
certify that no hidden correlations, drifts or external influences are present.

Quantum mechanics offers a fundamentally different source. According to the Born rule, the outcome of a measurement on a superposition state
is intrinsically unpredictable. No additional information about the system, even in principle, can determine it in advance (this statement is
made precise, and testable, by Bell's theorem, \cref{ch:bell}). This is the basis of \emph{quantum random number generators} (QRNGs) \cite{Jacak2021}.

\subsection{Example: a photon and a polarizer}
Consider a single photon linearly polarized at an angle $\theta$ with respect to the horizontal axis:
\begin{equation}
  \ket\psi=\cos\theta\ket H+\sin\theta\ket V.
\end{equation}
The photon is sent onto a polarizer that transmits horizontal polarization, followed by a single-photon detector. This realizes the
projective measurement with projectors $P_H=\proj H$ (transmitted, detector clicks) and $P_V=\proj V$ (absorbed, no click). The transmission
probability is
\begin{equation}
  p(H)=\mel\psi{P_H}\psi=|\braket H\psi|^2=\cos^2\theta,
\end{equation}
which is Malus's law at the level of a single photon. For $\theta=45^\circ$, $\ket\psi=(\ket H+\ket V)/\sqrt2$ and $p(H)=p(V)=1/2$. Each photon is
transmitted or absorbed with equal probability. Assigning the bit value $1$ to a click and $0$ to no click, a stream of identically prepared
photons produces a sequence of independent, unbiased random bits. After the measurement the transmitted photon is in state $\ket H$, so its
polarization no longer carries any memory of the random choice.

\subsection{Practical remarks}
A real device deviates from this ideal picture. If the angle is not exactly $45^\circ$, the bits are biased, with $p(1)=\cos^2\theta\neq\frac12$.
Photon losses and the finite detector efficiency also look like ``no click'' events and add to the bias. A common improvement replaces the
polarizer by a \emph{polarizing beam splitter} with a detector in each output port, so that both outcomes are actually detected and events with no
photon are discarded. Most commercial QRNGs actually use a beam splitter, which is the same measurement for a photon in a superposition of two
spatial modes.

Residual bias is removed by classical post-processing. The simplest method is \emph{von Neumann's procedure} \cite{vonNeumann1951}: group the bits
in pairs, map $01\to0$ and $10\to1$, and discard $00$ and $11$.
\begin{proposition}
For independent bits with any fixed bias $P(1)=p$, the output of the von Neumann extractor is exactly unbiased; on average it produces
$p(1-p)$ output bits per input bit.
\end{proposition}
\begin{proof}
For independent bits, $P(01)=(1-p)p=P(10)$, so conditioned on keeping the pair, both outputs have probability $1/2$. A pair is kept with
probability $2p(1-p)$ and yields one bit, i.e.\ $p(1-p)$ bits per input bit (at most $1/4$, for $p=1/2$).
\end{proof}
The method relies on independence; correlated bits require stronger \emph{randomness extractors} (e.g.\ hashing with Toeplitz matrices), whose
output length is set by the min-entropy $H_{\min}=-\log_2\max_xP(x)$ of the raw data.

Finally, the randomness of such a device rests on trusting its physical model (state preparation, polarizer and detectors). \emph{Device-independent}
QRNGs remove this assumption by certifying the randomness through the violation of a Bell inequality, at the cost of a much lower bit rate.

\section{The no-cloning theorem}\label{sec:nocloning}
A natural question is whether an unknown quantum state can be copied. Suppose there exists a linear operator $C$ that clones any state $\ket\alpha$
onto a second system initially prepared in a reference state $\ket0$:
\begin{equation}\label{eq:clone}
  C\ket\alpha\ket0=\ket\alpha\ket\alpha\qquad\text{for all }\ket\alpha.
\end{equation}
\begin{theorem}[No-cloning \cite{WoottersZurek1982,Dieks1982}]
No linear (in particular, no unitary) operator satisfies \cref{eq:clone} for all states.
\end{theorem}
\begin{proof}[Proof 1 (linearity)]
Let $\ket a$, $\ket b$ be two orthogonal states. By \cref{eq:clone}, $C\ket a\ket0=\ket a\ket a$ and $C\ket b\ket0=\ket b\ket b$. For the
superposition $\ket x=(\ket a+\ket b)/\sqrt2$, linearity gives
\begin{equation}
  C\ket x\ket0=\frac1{\sqrt2}\left(\ket a\ket a+\ket b\ket b\right),
\end{equation}
whereas a true clone would be
\begin{equation}
  \ket x\ket x=\frac12\left(\ket a\ket a+\ket a\ket b+\ket b\ket a+\ket b\ket b\right).
\end{equation}
The two differ: the first does not contain the cross terms, and even its coefficients differ. Contradiction.
\end{proof}
\begin{proof}[Proof 2 (unitarity)]
If $U\ket a\ket0=\ket a\ket a$ and $U\ket b\ket0=\ket b\ket b$, take the scalar product of the two equations. Since $U^\dagger U=I$, the left-hand side gives
$\braket ab\braket00=\braket ab$ and the right-hand side gives $\braket ab^2$. Thus $\braket ab=\braket ab^2$, so $\braket ab=0$ or $1$:
only identical or orthogonal states can be cloned by the same machine. (Allowing the machine to end in a state $\ket{M_a}$ that depends on the input does
not help: one gets $|\braket ab|=|\braket ab|^2|\braket{M_a}{M_b}|\le|\braket ab|^2$.)
\end{proof}
Cloning is therefore not forbidden for states known to belong to a fixed orthogonal set $\{\ket a,\ket b\}$: classical information can be copied (a CNOT
copies $\ket0$ and $\ket1$). It is the requirement to copy their superpositions as well that leads to the contradiction.

\begin{remark}[History and relation to relativity]
The theorem was found in 1982 in response to N.~Herbert's ``FLASH'' proposal for faster-than-light communication. Alice and Bob share an entangled pair
(\cref{ch:entanglement}); Alice measures her particle either in the $Z$ or in the $X$ basis, which leaves Bob's particle in a state of the $Z$ basis or of
the $X$ basis. If Bob could make many copies of his particle, he could measure them and find out which basis Alice chose, instantaneously. No-cloning
forbids it. More generally, perfect cloning would allow signalling and perfect state discrimination.
\end{remark}
\begin{remark}[Approximate cloning\starred]
Imperfect copies are allowed. The best universal cloner for qubits (Bužek and Hillery, 1996) produces two copies, each with fidelity
$F=\mel\psi{\rho_{\rm copy}}\psi=5/6$ for every input state. Its Bloch vector is shrunk by a factor $2/3$.
\end{remark}

\section{Quantum key distribution: BB84}\label{sec:bb84}
\subsection{The problem}
The only provably secure classical cipher is the \emph{one-time pad} (Vernam, 1917): the message $m$ is combined with a random key $k$ of the same length,
$c=m\oplus k$, and decrypted by $m=c\oplus k$. The ciphertext is statistically independent of the message (Shannon, 1949), provided that the key is
truly random, used only once, and \emph{known only to the two parties}. The hard part is to distribute the key. Public-key cryptography (RSA) solves
this problem only under computational assumptions, which Shor's algorithm (\cref{ch:shor}) would break.

The BB84 protocol, proposed by Bennett and Brassard in 1984 \cite{BB84}, was the first quantum key distribution (QKD) scheme. Its goal is to establish a
shared secret key between two parties, traditionally called Alice and Bob. Its security does not rely on computational assumptions. It follows from the
laws of quantum mechanics: an eavesdropper cannot copy an unknown quantum state in order to measure it later (no-cloning), and cannot measure it without
disturbing it if she does not know the basis.

The protocol uses two channels: a \emph{quantum channel}, along which Alice sends single photons to Bob, and a public but \emph{authenticated} classical
channel, along which they compare their choices of basis. (Authentication requires a short pre-shared secret; QKD is therefore really a key
\emph{expansion} protocol.) It was originally formulated for single photons. Ekert later proposed a version based on entangled pairs \cite{Ekert1991}
(\cref{sec:e91}).

\subsection{Encoding}
Alice encodes each bit in the polarization of a photon, using one of two bases:
\begin{align}
  \text{rectilinear basis }(+):&\quad 0\mapsto\ket\uparrow,\quad 1\mapsto\ket\rightarrow,\\
  \text{diagonal basis }(\times):&\quad 0\mapsto\ket\nearrow,\quad 1\mapsto\ket\searrow,
\end{align}
with $\ket\nearrow,\ket\searrow=(\ket\rightarrow\pm\ket\uparrow)/\sqrt2$. The states within one basis are orthogonal, while the two bases are
\emph{mutually unbiased}: $|\braket\uparrow\nearrow|^2=|\braket\uparrow\searrow|^2=1/2$. On the Poincaré sphere, the four states lie at the four points
$\pm\vec e_z$, $\pm\vec e_x$ of a great circle. If a photon is measured in the basis in which it was prepared, the result is certain. If it is measured in the
other basis, the result is completely random.

\subsection{The protocol}
\begin{enumerate}
  \item Alice generates a random bit string and, for each bit, a random choice of basis ($+$ or $\times$). She sends Bob one photon per bit, polarized accordingly.
  \item Bob independently chooses a random basis for each photon and measures its polarization in that basis.
  \item Over the public channel, Alice and Bob announce the bases they used, but not the bit values. They keep only the bits for which the bases coincide
  (on average half of them). This is the \emph{sifted key}.
  \item They publicly compare a random subset of the sifted key to estimate the \emph{quantum bit error rate} (QBER) $Q$, and discard these bits.
  \item If $Q$ is below a threshold, they perform classical \emph{error correction} (so that both have the same string) and \emph{privacy amplification}
  (hashing the string into a shorter one about which Eve's information is negligible).
\end{enumerate}
An example is shown in \cref{tab:bb84}. In columns 2, 4, 5 and 7 the bases differ, and Bob's result is random. In the other columns his result coincides with Alice's bit.

\begin{table}[h]\centering
\begin{tabular}{@{}lcccccccc@{}}\toprule
Alice's random bit & 0&1&1&0&1&0&0&1\\
Alice's basis & $+$&$+$&$\times$&$+$&$\times$&$\times$&$\times$&$+$\\
Polarization sent & $\uparrow$&$\rightarrow$&$\searrow$&$\uparrow$&$\searrow$&$\nearrow$&$\nearrow$&$\rightarrow$\\
Bob's basis & $+$&$\times$&$\times$&$\times$&$+$&$\times$&$+$&$+$\\
Polarization measured & $\uparrow$&$\nearrow$&$\searrow$&$\nearrow$&$\rightarrow$&$\nearrow$&$\rightarrow$&$\rightarrow$\\
Same basis? & \checkmark&--&\checkmark&--&--&\checkmark&--&\checkmark\\
Shared key & 0&&1&&&0&&1\\ \bottomrule
\end{tabular}
\caption{An example of the BB84 protocol with eight photons.}\label{tab:bb84}
\end{table}

\subsection{Security against intercept--resend}
Suppose an eavesdropper, Eve, intercepts each photon, measures it in a randomly chosen basis, and resends to Bob a photon in the state she found.
Consider a bit that survives sifting (Alice and Bob used the same basis).
\begin{itemize}
  \item With probability $1/2$ Eve chose the right basis: she learns the bit and resends the correct state. No error.
  \item With probability $1/2$ she chose the wrong basis: her result is random, and she resends a state of the wrong basis. Bob, measuring in Alice's basis,
  obtains a random result, wrong with probability $1/2$.
\end{itemize}
Eve therefore introduces errors in $\frac12\times\frac12=25\%$ of the sifted bits, while learning $50\%$ of them. Comparing $m$ bits, Alice and Bob
detect her with probability $1-(3/4)^m$, i.e.\ $>99.99\%$ for $m=32$. Because of the no-cloning theorem, Eve cannot avoid this disturbance by copying
the photons. More subtle attacks trade information for disturbance, but there is always a trade-off: \emph{information gain implies disturbance}.

\begin{caveat}[: how much error is tolerable?]
Real channels have errors even without Eve, and Alice and Bob must attribute \emph{all} errors to Eve. The Shor--Preskill security proof (2000)
shows that, against the most general attacks, a secret key can be distilled at the rate (secret bits per sifted bit)
\begin{equation}
  r=1-2h(Q),\qquad h(Q)=-Q\log_2Q-(1-Q)\log_2(1-Q),
\end{equation}
where one $h(Q)$ is the cost of error correction and the other the cost of privacy amplification. The rate vanishes at $Q\approx11\%$. Above this threshold
the key is discarded.
\end{caveat}

\subsection{Practical QKD}
Single-photon sources are difficult, so most systems use attenuated laser pulses, coherent states with a mean photon number $\mu\approx0.1$--$0.5$. The number of
photons per pulse is Poisson distributed, $P(n)=\e^{-\mu}\mu^n/n!$, so a fraction $P(n\ge2)/P(n\ge1)\approx\mu/2$ of the non-empty pulses contains several photons.
Eve can keep one of them and let the others go to Bob (\emph{photon-number-splitting attack}). The \emph{decoy-state} method (Hwang 2003; Lo, Ma, Chen 2005)
defeats this attack by randomly varying $\mu$, which lets Alice and Bob estimate the single-photon contribution. Detector side channels (``blinding'' attacks)
motivate measurement-device-independent QKD. Commercial QKD systems reach a few hundred kilometres in fibre; the Chinese satellite Micius distributed keys
over $1200$\,km (2017). Beyond that, \emph{quantum repeaters}, based on entanglement swapping (\cref{ch:teleportation}) and quantum memories, are required, since
the no-cloning theorem also forbids optical amplification of quantum signals.

\section{Exercises}
\begin{exercise}
Show that a half-wave plate at $22.5^\circ$ implements the Hadamard gate on polarization qubits, and that a QWP at $45^\circ$ maps $\ket H$ to a circular polarization.
\end{exercise}
\begin{exercise}
In BB84, Eve measures every photon in the \emph{Breidbart basis}, halfway between the two bases (at $22.5^\circ$ in real space). Show that she guesses
each sifted bit correctly with probability $\cos^2(\pi/8)\approx0.85$, and that she induces a QBER of $25\%$.
\end{exercise}
\begin{exercise}
Compute the probability that Eve, measuring every photon as in the intercept--resend attack, remains undetected when Alice and Bob compare $m$ bits. How many bits must
be compared to detect her with probability $1-10^{-9}$?
\end{exercise}
